% !TeX program = xelatex
% ============================================================================
%  第 8 课　正方形对称：D₄ 的 8 个元素　—— 学生版讲义
%  与 02-课件/第08课-正方形的对称D4/第08课-正方形的对称D4.tex 同步
% ============================================================================

\zhipianye{8}

\xhlesson{8}{正方形的对称：\texorpdfstring{\(D_4\)}{D4} 的 \texorpdfstring{\(8\)}{8} 个元素}{}

\begin{mubiao}
  \begin{itemize}
    \item 能说出正方形的 \(8\) 个对称动作。
    \item 会把这 \(8\) 个动作写成两行式。
    \item 知道做法和三角形完全一样。
  \end{itemize}
\end{mubiao}

\section*{一、八个动作}

\noindent 正方形四个顶点写 \(1,2,3,4\)：\quad \(1\) 左上，\(2\) 右上，\(3\) 右下，\(4\) 左下。

% 待插入 \imgfig[0.5\textwidth]{zhipian_zhengfangxing}：一张正方形纸片示意（正面 1 左上、2 右上、3 右下、4 左下，背面另一种颜色）
\begin{dongshou}[边看边记]
  老师每做一个动作就贴回原位。\\
  转动那 \(4\) 行跟着老师做、自己填；翻折那 \(4\) 行先空着，等老师翻开这一页\textbf{对答案}。
  四个顶点写不下就写简写（比如「\(2\)」表示 \(1\to 2\)）。

  \begin{center}
  {\small
  \begin{tabular}{@{}lc@{}}
    \toprule
    动作 & 顶点去了哪（只记 \(1\) 和 \(3\)） \\
    \midrule
    \rule{0pt}{2.6ex}不动 & \(1\to\underline{\hspace{1.2em}}\)，\(3\to\underline{\hspace{1.2em}}\) \\[3pt]
    \rule{0pt}{2.6ex}顺时针转 \(90^\circ\) & \(1\to\underline{\hspace{1.2em}}\)，\(3\to\underline{\hspace{1.2em}}\) \\[3pt]
    \rule{0pt}{2.6ex}顺时针转 \(180^\circ\) & \(1\to\underline{\hspace{1.2em}}\)，\(3\to\underline{\hspace{1.2em}}\) \\[3pt]
    \rule{0pt}{2.6ex}顺时针转 \(270^\circ\) & \(1\to\underline{\hspace{1.2em}}\)，\(3\to\underline{\hspace{1.2em}}\) \\
    \midrule
    \rule{0pt}{2.6ex}沿对角线 \(1\text{—}3\) 翻（对答案） & \(1\to 1\)，\(3\to 3\) \\[3pt]
    \rule{0pt}{2.6ex}沿对角线 \(2\text{—}4\) 翻（对答案） & \(1\to 3\)，\(3\to 1\) \\[3pt]
    \rule{0pt}{2.6ex}沿竖直中线翻（对答案） & \(1\to 2\)，\(3\to 4\) \\[3pt]
    \rule{0pt}{2.6ex}沿水平中线翻（对答案） & \(1\to 4\)，\(3\to 2\) \\
    \bottomrule
  \end{tabular}}
  \end{center}
  自查：\(1\)、\(3\) 的去向都记对了，另外两个顶点也就跟着定下来了。
\end{dongshou}

\begin{sikao}[反例]
  把正方形转过 \(45^\circ\)，为什么不算一个对称动作？\\
  （想一想第 5 课三角形「歪着挪一下」为什么也不算。）
  \liukong{2}
\end{sikao}

\begin{sikao}
  一边数一边答：正方形一共有 \underline{\hspace{1.6em}} 个动作；
  其中转动 \underline{\hspace{1.2em}} 个，翻折 \underline{\hspace{1.2em}} 个。\\[4pt]
  转 \(90^\circ\) 要连做几次才回到原位？\underline{\hspace{1.6em}} 次。
  三角形里转 \(120^\circ\) 呢？\underline{\hspace{1.6em}} 次。
\end{sikao}

\section*{二、八张两行式}

\begin{example}
  取名字：\(r\) 是「顺时针转 \(90^\circ\)」，\(f\) 是「沿对角线 \(1\text{—}3\) 翻折」。
  \begin{center}
  {\small
  \begin{tabular}{@{}lcc@{}}
    \toprule
    名字 & 动作 & 两行式 \\
    \midrule
    \(e\) & 不动 & \(\begin{pmatrix}1&2&3&4\\1&2&3&4\end{pmatrix}\) \\[5pt]
    \(r\) & 转 \(90^\circ\) & \(\begin{pmatrix}1&2&3&4\\2&3&4&1\end{pmatrix}\) \\[5pt]
    \(r^2\) & 转 \(180^\circ\) & \(\begin{pmatrix}1&2&3&4\\3&4&1&2\end{pmatrix}\) \\[5pt]
    \(r^3\) & 转 \(270^\circ\) & \(\begin{pmatrix}1&2&3&4\\4&1&2&3\end{pmatrix}\) \\[5pt]
    \(f\) & 沿对角线 \(1\text{—}3\) 翻 & \(\begin{pmatrix}1&2&3&4\\1&4&3&2\end{pmatrix}\) \\[5pt]
    \(rf\) & 沿水平中线翻 & \(\begin{pmatrix}1&2&3&4\\4&3&2&1\end{pmatrix}\) \\[5pt]
    \(r^2f\) & 沿对角线 \(2\text{—}4\) 翻 & \(\begin{pmatrix}1&2&3&4\\3&2&1&4\end{pmatrix}\) \\[5pt]
    \(r^3f\) & 沿竖直中线翻 & \(\begin{pmatrix}1&2&3&4\\2&1&4&3\end{pmatrix}\) \\
    \bottomrule
  \end{tabular}}
  \end{center}
  顺手说一句：\(rf\)、\(r^2f\)、\(r^3f\) 这\textbf{三行可以自己用两行式推出来}——
  思路和三角形的 \(rf\)、\(r^2f\) 一样，一步一步跟着顶点走就行。
\end{example}

\section*{三、填乘法表}

\begin{dongshou}[只填 \(r\) 行与 \(f\) 行]
  格子里的动作是「先横排的，再做竖排的」。

  \begin{center}
  {\small
  \begin{tabular}{@{}l|cccccccc@{}}
    \toprule
     & \(e\) & \(r\) & \(r^2\) & \(r^3\) & \(f\) & \(rf\) & \(r^2f\) & \(r^3f\) \\
    \midrule
    \rule{0pt}{2.8ex}先 \(r\) & & & & & & & & \\
    \midrule
    \rule{0pt}{2.8ex}先 \(f\) & & & & & & & & \\
    \midrule
    \rule{0pt}{2.8ex}先 \(r^2\)（挑战） & & & & & & & & \\
    \bottomrule
  \end{tabular}}
  \end{center}
  前两行是必做；\textbf{第三行「先 \(r^2\)」是挑战行}，做得动就填，做不动就空着。
  （提示：先 \(r^2\) 就是「先转 \(180^\circ\)」。\(r^2\) 再接着做一次 \(r^2\)，等于转 \(360^\circ\)，就回到原样了。）\\
  一列一格地写，别串行。\textbf{列窄，写不下就写简写}。\\
  （比如 \(r^3f\) 写不下时，就写成「\(3f\)」，自己看得懂就行。）
\end{dongshou}

\begin{sikao}
  \(r\) 那一行里，有没有哪个动作出现了两次？\underline{\hspace{3em}}\\
  所以这一行是 \(8\) 个动作的\underline{\hspace{6em}}。
\end{sikao}

\begin{example}
  从 \(f\) 行的第二个格子读出一条关系：
  \[
    fr=r^3f .
  \]
  和三角形的 \(fr=r^2f\) 比一比：两边都是「先翻后转，等于先\textbf{倒着}转再翻」。
\end{example}

\begin{xiaojie}
  列一张小对照表：三角形的动作群与正方形的动作群，各有多少个动作？基本转动各转多少度？
  \liukong{3}
\end{xiaojie}

\noindent\textbf{下节课}：在正方形里找小群（子群），再看旋转与翻折那条关系。